\section*{Capítulo \ref{ch:b1:flowering-plant-organization} --- Organização funcional de uma planta com flores}

\begin{solution}{exo:b1:flowering-plant-organization:1}
\omterm{def:b1:flowering-plant-organization:organs}{Raiz}: fixação, absorção de água e de íons minerais. \omterm{def:b1:flowering-plant-organization:organs}{Caule}: sustentação
das \omterm{def:b1:flowering-plant-organization:organs}{folhas} em direção à luz, condução entre raízes e \omterm{def:b1:flowering-plant-organization:organs}{folhas}. \omterm{def:b1:flowering-plant-organization:organs}{Folha}:
fotossíntese (e trocas gasosas, transpiração).
\end{solution}

\begin{solution}{exo:b1:flowering-plant-organization:2}
Um módulo é um \omterm{def:b1:flowering-plant-organization:organs}{nó} com a \omterm{def:b1:flowering-plant-organization:organs}{folha} e a gema axilar dele, mais o \omterm{def:b1:flowering-plant-organization:organs}{entrenó} de
baixo. \omterm{prop:b1:flowering-plant-organization:modular}{Crescimento indeterminado}: não há tamanho adulto; as gemas apical
e axilares continuam acrescentando módulos enquanto a planta viver.
\end{solution}

\begin{solution}{exo:b1:flowering-plant-organization:3}
Epiderme: viva; revestimento, \omterm{def:b1:flowering-plant-organization:tissues}{cutícula}, \omterm{def:b1:flowering-plant-organization:tissues}{estômatos}, \omterm{def:b1:flowering-plant-organization:organs}{pelos radiculares}.
\omterm{def:b1:flowering-plant-organization:tissues}{Parênquima}: vivo; fotossíntese, armazenamento. \omterm{def:b1:flowering-plant-organization:tissues}{Colênquima}: vivo;
sustentação flexível. \omterm{def:b1:flowering-plant-organization:tissues}{Esclerênquima}: morto; sustentação rígida. Células
condutoras do \omterm{def:b1:flowering-plant-organization:tissues}{xilema}: mortas; transporte de água. \omterm{def:b1:flowering-plant-organization:tissues}{Tubos crivados} do
\omterm{def:b1:flowering-plant-organization:tissues}{floema}: vivos (com as células companheiras); transporte de açúcar.
\end{solution}

\begin{solution}{exo:b1:flowering-plant-organization:4}
$e^{-0.9} = 0.41$, ou seja, \qty{41}{\%}; $e^{-3} = 0.05$, ou seja,
\qty{5}{\%}.
\end{solution}

\begin{solution}{exo:b1:flowering-plant-organization:5}
Uma \omterm{def:b1:flowering-plant-organization:organs}{raiz} (cilindro central único, \omterm{prop:b1:flowering-plant-organization:stemroot}{endoderme}, \omterm{def:b1:flowering-plant-organization:tissues}{xilema} e \omterm{def:b1:flowering-plant-organization:tissues}{floema} alternados
sobre um raio) de dicotiledônea (poucos braços de \omterm{def:b1:flowering-plant-organization:tissues}{xilema}, sem medula
central).
\end{solution}

\begin{solution}{exo:b1:flowering-plant-organization:6}
Uma viga fletida é tensionada ao máximo na superfície e nada ao longo do
eixo, de modo que material colocado num anel periférico resiste à flexão
com a menor massa possível (um tubo). Um cabo sob tração é tensionado
uniformemente em toda a seção e se dobra livremente; um cordão central
suporta a tração enquanto o córtex permanece flexível, e as raízes
laterais podem emergir através dele.
\end{solution}

\begin{solution}{exo:b1:flowering-plant-organization:7}
$1 - e^{-2.4} = 0.91$; com $L = 6$, $1 - e^{-3.6} = 0.97$: duas camadas
de \omterm{def:b1:flowering-plant-organization:organs}{folha} a mais acrescentam \qty{6}{\%} da luz e custam \qty{50}{\%}
mais \omterm{def:b1:flowering-plant-organization:organs}{folha}.
\end{solution}

\begin{solution}{exo:b1:flowering-plant-organization:8}
$S = \pi d L$: $d = 240/(\pi\times 623\,000) = \qty{1.2e-4}{m}$, cerca
de \qty{0.12}{mm}. $(240 + 400)/5 = 128$.
\end{solution}

\begin{solution}{exo:b1:flowering-plant-organization:9}
Superfície de \omterm{def:b1:flowering-plant-organization:organs}{raiz} por mm: $\pi\times 0.4 = \qty{1.26}{mm^2}$. Pelos por
mm: $250\times\pi\times 0.010\times 0.8 = \qty{6.3}{mm^2}$. Total
$7.5/1.26 = 6$: os pelos multiplicam a superfície por seis.
\end{solution}

\begin{solution}{exo:b1:flowering-plant-organization:10}
Na floresta as gemas axilares inferiores ficam sombreadas; os módulos
que elas produziriam não conseguem se pagar, de modo que elas
permanecem dormentes ou os ramos delas morrem, e o crescimento é
destinado à gema apical, que corre para cima em direção à luz. No campo
toda gema é iluminada e todo ramo se paga, de modo que a copa se
preenche desde baixo. O custo: um tronco alto e fino é um grande
investimento em madeira improdutiva, vulnerável ao vento; o ganho: a
planta ajusta a forma dela a onde está a luz, sem nenhum plano fixo.
\end{solution}

\begin{solution}{exo:b1:flowering-plant-organization:11}
A cripta e os pelos retêm uma camada de ar parado e úmido acima do poro;
o gradiente de vapor entre os espaços de ar e o ar em movimento se
distribui por um caminho mais longo e a transpiração cai. O dióxido de
carbono precisa se difundir para dentro pelo mesmo caminho mais longo,
de modo que a captação também cai, embora menos, já que o gradiente de
$\mathrm{CO_2}$ é fixado pela atmosfera e pelo consumo da \omterm{def:b1:flowering-plant-organization:organs}{folha}. Num
deserto a água é o recurso limitante e a troca salva a vida da planta;
numa floresta tropical úmida a água é gratuita e é o ganho de carbono
que a competição recompensa.
\end{solution}

\begin{solution}{exo:b1:flowering-plant-organization:12}
A planta não tem um \omterm{def:b1:mammal-organization:compartments}{líquido extracelular} regulado, mantido a temperatura
e composição constantes entre as células dela e o exterior: as células
dela encontram diretamente a água do solo, o ar e a luz do sol, e a
temperatura dela é a do ar. Cada célula, portanto, regula o próprio
conteúdo (íons, água, pH) através da membrana e da parede; a planta
regula as trocas na superfície (os \omterm{def:b1:flowering-plant-organization:tissues}{estômatos} abrem e fecham, a \omterm{prop:b1:flowering-plant-organization:stemroot}{endoderme}
seleciona o que entra no \omterm{def:b1:flowering-plant-organization:tissues}{xilema}) e ajusta o corpo por crescimento (mais
raízes em solo seco, mais \omterm{def:b1:flowering-plant-organization:organs}{folhas} na sombra). Ela regula, mas a regulação
está nas superfícies e na forma, e não num líquido particular.
\end{solution}

\begin{solution}{pb:b1:flowering-plant-organization:1}
\textbf{1.} $5\times 100 = \qty{500}{m^2}$.
\textbf{2.} $e^{-2.5} = 0.082$ chegam ao solo; \qty{92}{\%}
interceptados.
\textbf{3.} $L = 2$: $1 - e^{-1} = \qty{63}{\%}$; $L = 8$: $1 -
e^{-4} = \qty{98}{\%}$.
\textbf{4.} De 2 a 5 a copa ganha \qty{29}{\%} da luz por
\qty{300}{m^2} de \omterm{def:b1:flowering-plant-organization:organs}{folha}; de 5 a 8 ela ganharia \qty{6}{\%} por mais
\qty{300}{m^2}, que custam o mesmo para construir e para manter
abastecidos de água: as \omterm{def:b1:flowering-plant-organization:organs}{folhas} a mais não se pagam.
\textbf{5.} $500/0.005 = \num{100000}$ \omterm{def:b1:flowering-plant-organization:organs}{folhas}.
\textbf{6.} $500\times 10^6\,\mathrm{mm^2}\times 150 = \num{7.5e10}$
\omterm{def:b1:flowering-plant-organization:tissues}{estômatos}.
\textbf{7.} $\qty{50}{\micro m^2}$; por mm$^2$: $150\times 50\times
10^{-6} = \qty{7.5e-3}{mm^2}$, ou seja, \qty{0.75}{\%} da superfície.
\textbf{8.} $500\times 4\times 10^{-6} = \qty{2e-3}{mol/s}$, ou seja,
\qty{88}{mg/s}.
\textbf{9.} $0.088\times 43\,200 = \qty{3.8}{kg}$ de $\mathrm{CO_2}$,
contendo $3.8\times 12/44 = \qty{1.04}{kg}$ de carbono.
\textbf{10.} $150\times 1.04 = \qty{156}{kg}$ de carbono; a metade,
\qty{78}{kg}, em madeira com \qty{50}{\%} de carbono: \qty{156}{kg} de
madeira.
\textbf{11.} $500\times 10^{-3} = \qty{0.5}{mol/s}$, ou seja, \qty{9}{g/s}.
\textbf{12.} $9\times 43\,200 = \qty{389}{kg}$, cerca de \qty{390}{L}.
\textbf{13.} $0.5/(2\times 10^{-3}) = 250$ moléculas de água por
$\mathrm{CO_2}$. O poro que deixa o $\mathrm{CO_2}$ entrar deixa a água
sair; o ar de dentro está saturado, ao passo que o $\mathrm{CO_2}$ é
apenas \qty{0.04}{\%} do ar de fora, de modo que o gradiente de vapor
para fora é muito mais íngreme que o gradiente de $\mathrm{CO_2}$ para
dentro.
\textbf{14.} $\qty{2}{mm}\times\qty{100}{m^2} = \qty{200}{L}$: meio dia.
\textbf{15.} \qty{100}{m^3}; $5\times 100 = \qty{500}{km}$.
\textbf{16.} $\pi\times 0.5\times 10^{-3}\times 5\times 10^5 =
\qty{785}{m^2}$.
\textbf{17.} $200\times 5\times 10^8\,\mathrm{mm} = \num{1e11}$ pelos.
\textbf{18.} Cada um $\pi\times 10^{-5}\times 5\times 10^{-4} =
\qty{1.57e-8}{m^2}$; total \qty{1570}{m^2}.
\textbf{19.} $785 + 1570 = \qty{2355}{m^2}$, $4.7$ vezes a área foliar.
\textbf{20.} Volume a menos de \qty{1}{mm}: $\pi\times(10^{-3})^2\times 5
\times 10^5 = \qty{1.6}{m^3}$, ou seja, \qty{1.6}{\%} do solo.
\textbf{21.} Centeio: $\pi\times 10^{-6}\times 6.2\times 10^5 =
\qty{1.95}{m^3}$ do vaso de \qty{0.05}{m^3}, ou seja, \qty{3900}{\%}:
cada ponto do solo está a menos de um milímetro de quarenta raízes. O
centeio explora o solo dele mais de duas mil vezes mais a fundo --- uma
planta de cultura com um orçamento de quatro meses contra uma árvore que
explora cem metros cúbicos.
\textbf{22.} $1000 + 2355 = \qty{3355}{m^2}$; $\qty{34}{m^2}$ por metro
quadrado de solo.
\textbf{23.} Ser humano: $130/2 = 65$; a razão do carvalho é a metade da
das superfícies dobradas do ser humano, sem nenhum dobramento.
\textbf{24.} O carvalho acrescenta módulos, \omterm{def:b1:flowering-plant-organization:organs}{folhas} e raízes todo ano, de
modo que as superfícies dele crescem com a idade; as superfícies do
mamífero são \omterm{def:b1:mammal-organization:organ}{órgãos} fixos, completos no tamanho adulto, e a única
maneira de elevar a razão dele foi dobrá-las durante o desenvolvimento.
\textbf{25.} Cerca de \qty{34}{m^2} de \omterm{prop:b1:organism-environment:surfaces}{superfície de troca} por metro
quadrado de solo: \qty{10}{m^2} de \omterm{def:b1:flowering-plant-organization:organs}{folha} (as duas faces) acima e
\qty{24}{m^2} de \omterm{def:b1:flowering-plant-organization:organs}{raiz} e \omterm{def:b1:flowering-plant-organization:organs}{pelo radicular} abaixo.
\end{solution}
